\documentclass[11pt,a4paper]{article}
\usepackage[margin=2.5cm]{geometry}
\usepackage{amsmath,amssymb,amsthm}
\usepackage{booktabs}
\usepackage{graphicx}
\usepackage{xcolor}
\usepackage[hidelinks]{hyperref}
\usepackage{url}

\newtheorem{lemma}{Lemma}
\newtheorem{proposition}{Proposition}
\newtheorem{corollary}{Corollary}
\theoremstyle{definition}
\newtheorem{definition}{Definition}
\theoremstyle{remark}
\newtheorem{remark}{Remark}

\newcommand{\add}{\boxplus}
\newcommand{\F}{\mathbb{F}_2}
\newcommand{\val}{\operatorname{val}}
\newcommand{\repo}{\url{https://github.com/AlexGuseinov/dipper-integral}}

\title{Certified Integral Properties of the Dipper Block Cipher:\\
Monomial-Trail Analysis and the Role of the Half-State Modular Addition}

% Author list to be confirmed by the authors.
\author{Ali Huseynli\thanks{Corresponding author. ORCID 0009-0003-4959-2887.}
\and Yadigar Imamverdiyev \and Jalal Alizadeh\\[4pt]
\small Department of Cybersecurity, Faculty of Information Technologies and Telecommunications,\\
\small Azerbaijan Technical University, Baku AZ1073, Azerbaijan}
\date{}

\begin{document}
\maketitle

\begin{abstract}
Dipper is a 64-bit lightweight block cipher whose round combines the GIFT S-box and bit permutation with two 16-bit modular additions acting on half of the state. In the design paper we bounded integral distinguishers experimentally at five rounds. We replace this experimental evidence with machine-checked certificates that hold for every choice of the round keys. We model monomial trails through reduced-round Dipper as SAT instances, using the exact coefficient table of the S-box and an exact local rule for the map $(x,y)\mapsto(x\add y,\,y)$ that we derive from the Braeken--Semaev/Hu--Yap characterisation of modular addition and verify exhaustively. The models are validated component-wise against exact algebraic normal forms, against published division-property results for GIFT-64 and PRESENT, and against cube-sum experiments with random keys, which never contradicted a certificate. On Dipper we obtain: (i) the published five-round properties are reproduced and certified for all keys; (ii) six-round integral properties exist (up to nine balanced bits with $2^{63}$ chosen plaintexts, six balanced bits with $2^{60}$); (iii) no seven-round property is certified for any bit-aligned cube and single output bit, which is exhaustive within the model because of a monotonicity lemma; (iv) since Dipper has no final key whitening, every $r$-round property extends to $r{+}1$ rounds without key guessing, giving a seven-round distinguisher and an eight-round partial key recovery with $2^{63}$ data. In a controlled ablation, replacing the two modular additions by XOR extends the longest certified property to nine rounds, and removing them extends it to ten; all certified balanced bits at the six-round boundary originate from the two words that bypass the addition in the last round. Exact monomial prediction and the conventional bit-based division property produce identical certificates in every case we examined, which we partly explain by the full-state key addition that precedes each S-box layer. We also report that the published Dipper-64/96 test vectors could not be reproduced from the specification text and that the 128-bit vectors require a different round-key window than the one stated. All code and data are public.
\end{abstract}

\noindent\textbf{Keywords:} integral cryptanalysis; division property; monomial prediction; modular addition; SAT; lightweight block cipher; Dipper

\section{Introduction}

Integral (square, saturation) cryptanalysis~\cite{knudsen2002integral} looks for a set of chosen plaintexts, typically an affine subspace (a \emph{cube}), over which some output bit sums to zero independently of the key. The division property~\cite{todo2015} and its bit-based versions~\cite{todo2016bit,xiang2016milp} turned the search for such properties into a propagation problem that can be solved with MILP or SAT tools; monomial prediction~\cite{hu2020mp}, shown in~\cite{hu2020mp} to be equivalent to the three-subset division property without unknown subset~\cite{hao2020}, gives the exact algebraic view in which a cube sum vanishes if and only if the corresponding monomial coefficient is zero.

Dipper, proposed in our previous work~\cite{dipper2026}, is a 64-bit hybrid SPN--ARX cipher with 96- and 128-bit keys and 28 rounds. Each round applies a full-state key addition, sixteen GIFT S-boxes~\cite{gift2017}, four word rotations, two 16-bit modular additions over half of the state and the GIFT-64 bit permutation. Its security evaluation includes MILP differential bounds, CP-SAT linear bounds, an impossible-differential search and an \emph{experimental} integral search that saturates one 16-bit word and reports distinguishers of at most five rounds. Experimental integral searches have two weaknesses: a sum that is zero for a few random keys is not necessarily zero for all keys, and a search limited to word-saturated cubes says nothing about other cubes.

\paragraph{Contributions.}
\begin{enumerate}
\item \textbf{Exact local models.} We model the Dipper round with the exact $16\times16$ monomial table of the GIFT S-box and an exact local rule (Lemma~\ref{lem:addret}) for the addition with retained operand $(x,y)\mapsto(x\add y,y)$, which is the actual shape of Dipper's ARX layer. Every local model is checked exhaustively against algebraic normal forms.
\item \textbf{All-key certificates and a complete round bound within the model.} We certify the published five-round properties for all keys, find six-round properties, and prove that no bit-aligned cube of any dimension gives a certificate for a single output bit at seven rounds in our model (Lemma~\ref{lem:mono} reduces this to the 64 cubes of dimension~63).
\item \textbf{Free final round and key recovery.} Because Dipper has no final whitening key, every $r$-round state property is an $(r{+}1)$-round property of $T^{-1}(C)$ without key guessing (Lemma~\ref{lem:free}). We quantify the resulting seven-round distinguisher and an eight-round partial key recovery, and we validate the key-recovery procedure end-to-end at small scale.
\item \textbf{Model comparison and tightness.} Exact monomial prediction, gate-level monomial prediction and the conventional bit-based division property give identical certificates on every instance examined. We explain this and measure how close the certificates are to experiment.
\item \textbf{Role of the modular addition.} An ablation with the additions replaced by XOR or removed isolates the contribution of the modular additions: 6 certified rounds versus 9 and 10.
\item \textbf{Specification check.} We report discrepancies between the Dipper specification text and its published test vectors (Section~\ref{sec:erratum}).
\end{enumerate}
The repository \repo{} contains the reference implementation, all models, scripts and raw results.

\paragraph{Scope of the claims.} A certificate proves that a cube sum is zero for all keys. The absence of a certificate proves nothing about the cipher: it only means that the model found no certificate. Our round bound at seven rounds is therefore a statement about the models used here, not a proof that Dipper has no seven-round integral distinguisher.

\section{Preliminaries}

\subsection{Dipper}
The state is $S=A\|B\|C\|D$ with 16-bit words, $A=S[63{:}48]$ and bit~0 the least significant. For $r=1,\dots,28$,
\[
S^{(r)} = T\big(S^{(r-1)}\oplus RK_r\big),\qquad
T = P\circ M\circ \mathrm{SC},
\]
where $\mathrm{SC}$ applies the GIFT S-box to the sixteen nibbles, $M$ rotates $A,B,C,D$ left by $1,4,7,11$ and then sets $A\gets A\add B$ and $C\gets C\add D$ (with $B,D$ unchanged), and $P$ is the GIFT-64 bit permutation. The ciphertext is $S^{(28)}$; there is no final key addition. $T$ is a public permutation. We write $\mathrm{Dipper}^{\oplus}$ for the variant in which the two additions are replaced by XOR and $\mathrm{Dipper}^{\varnothing}$ for the variant in which they are removed (rotations kept). These are analysis variants only.

\subsection{Cube sums and monomial prediction}
For $u\in\F^n$ write $x^u=\prod_{i:u_i=1}x_i$. Let $I\subseteq\{0,\dots,63\}$ be a set of active bits and let the remaining plaintext bits be constants. For an output bit $f$ of $r$ rounds, viewed as a polynomial in the plaintext $x$ and the key material $k$,
\[
\bigoplus_{x_I\in\F^{|I|}} f(x,k) \;=\; \sum_{u\supseteq I,\,v} a_{u,v}\, x^{u\setminus I}k^v ,
\]
so the sum vanishes for all constants and all keys if and only if $a_{u,v}=0$ whenever $u\supseteq I$. For a composite function $f=f_r\circ\dots\circ f_1$, the coefficient of $x^{u_0}$ in $f^{u_r}$ equals the parity of the number of \emph{monomial trails} $u_0\to u_1\to\dots\to u_r$ in which each step is a nonzero local coefficient~\cite{hu2020mp}. If no trail exists, the coefficient is zero. This is the only direction we use: \emph{no trail} $\Rightarrow$ \emph{balanced}.

The conventional bit-based division property (BDP)~\cite{todo2016bit,xiang2016milp} propagates a set of vectors instead of monomials. It is sound in the same direction (no division trail to a unit vector $\Rightarrow$ the bit is balanced) and can be less precise.

\section{Models}

\subsection{Local propagation rules}

\paragraph{Key addition.} For $s'=s\oplus k$ with free $k$, $s'^{w}=\sum_{u\preceq w}s^u k^{w\oplus u}$. Hence $u\to w$ is a trail step iff $u\preceq w$, and distinct steps carry distinct key monomials. Under BDP the key addition is transparent.

\paragraph{S-box.} For $y=S(x)$ we compute $T[u][v]=[x^u]\,y^v$ for all $u,v\in\F^4$ by Möbius transform. For the GIFT S-box, 68 of the 256 pairs are nonzero. The BDP table is $\mathrm{BDP}[k][v]=1$ iff $T[u][v]=1$ for some $u\succeq k$, which gives 170 transitions.

\paragraph{Modular addition.} Braeken and Semaev~\cite{braeken2005}, applied to cryptanalysis by Hu and Yap~\cite{huyap2024}, show that for $z=x\add y$ on $n$ bits, $[x^ay^b]\,z^w=1$ if and only if $\val(a)+\val(b)=\val(w)$ over the integers. Dipper does not use $z$ alone: the addend is also an output. The following lemma gives the exact rule for that shape.

\begin{lemma}\label{lem:addret}
Let $F(x,y)=(x\add y,\,y)$ on $n$-bit words. Then $[x^a y^b]\,(x\add y)^w y^c = 1$ if and only if $\val(w)\ge\val(a)$ and $b = b'\vee c$, where $b'$ is the word with $\val(b')=\val(w)-\val(a)$.
\end{lemma}
\begin{proof}
By the Braeken--Semaev characterisation, $(x\add y)^w=\sum_{\val(a')+\val(b')=\val(w)}x^{a'}y^{b'}$. Multiplying by $y^c$ and using $y_i^2=y_i$ gives $\sum x^{a'}y^{b'\vee c}$. The coefficient of $x^ay^b$ is the parity of $\#\{b' : \val(a)+\val(b')=\val(w),\ b'\vee c=b\}$. The equation determines $b'$ uniquely, and $b'$ exists iff $\val(w)\ge\val(a)$.
\end{proof}
In the SAT model, the integer equation $\val(a)+\val(b')=\val(w)$ is encoded bitwise with auxiliary carries: $a_i+b'_i+q_i=w_i+2q_{i+1}$ with $q_0=q_n=0$. These are carries between exponent integers, not the cipher's data carries. Since $b'$ and $q$ are determined by $(a,w)$, each trail has exactly one extension to the auxiliary variables.

We call the model built from these tables \emph{exact MP}: its local relations are exact, while the global test is trail \emph{existence}, not trail parity.

\paragraph{Gate-level alternatives.} For comparison we also model the addition as a ripple-carry circuit ($s_i=x_i\oplus y_i\oplus c_i$, $c_{i+1}=x_iy_i\oplus(x_i\oplus y_i)c_i$) with explicit COPY gates, including the copy of $y$ that forms the retained output. The circuit is propagated with monomial rules (COPY: $u=\bigvee v_j$; AND: $u_1=u_2=v$; XOR: $v=u_1+u_2$) in the model we call \emph{MP-circuit}, and with BDP rules (COPY: $u=\sum v_j$; AND: $v=u_1\vee u_2$) in the model we call \emph{BDP}.

\subsection{Round model, soundness and two structural lemmas}
One round is modelled as key addition, sixteen S-box relations, the word rotations (wiring), the two addition relations and the bit permutation (wiring). The input mask is fixed to the indicator of $I$ and the output mask after $r$ rounds to the unit vector $e_j$. Every relation is a set of clauses that forbid the invalid assignments. We solve the instances with CaDiCaL through PySAT~\cite{pysat2018}. Each output bit is an incremental query under assumptions.

\begin{proposition}[Soundness]
If the exact MP model (or the BDP model) is unsatisfiable for $(I,r,j)$, then $\bigoplus_{x_I}S^{(r)}_j=0$ for every value of the constant bits and every sequence of round keys $RK_1,\dots,RK_r$. In particular, the property holds under both Dipper key schedules.
\end{proposition}
\begin{proof}
Every local relation contains all transitions with nonzero coefficient, so every nonzero coefficient of $x^{u}$, $u\supseteq I$, in $S_j^{(r)}$ has at least one trail starting from $u$. Such a trail first passes the key addition to some $w\succeq u\succeq I$, and $I\preceq w$ is also a valid key-addition step, so the model, whose input mask is exactly $I$, contains a trail as well. (Under BDP the same holds because the division property is defined by $u\succeq I$.) The round keys are independent free variables, so the statement covers all key sequences, and therefore the subset produced by any key schedule. The constant bits merge with $RK_1$ into free variables.
\end{proof}

\begin{lemma}[Monotonicity]\label{lem:mono}
In the MP model, if $I\subseteq I'$ and $(I,r,j)$ has no trail, then $(I',r,j)$ has no trail. Consequently, if no bit-aligned cube of dimension 63 yields a certificate for output bit $j$ after $r$ rounds, no bit-aligned cube of any dimension $\le63$ does.
\end{lemma}
\begin{proof}
The first operation is the key addition, so the masks reachable after it from $I$ are $\{w\succeq I\}$, which contains $\{w\succeq I'\}$. All subsequent constraints are identical. Every cube of dimension $\le 63$ lies inside some cube of dimension~63.
\end{proof}
The 64-dimensional cube (the full codebook) is trivially balanced for any permutation and is excluded.

\begin{lemma}[Free final round]\label{lem:free}
If $\bigoplus_{x\in\mathcal C}S^{(r)}_j=0$ for all keys and $|\mathcal C|$ is even, then $\bigoplus_{x\in\mathcal C}T^{-1}(S^{(r+1)})_j=0$ for all keys.
\end{lemma}
\begin{proof}
$T^{-1}(S^{(r+1)})=S^{(r)}\oplus RK_{r+1}$, and the constant $RK_{r+1}$ cancels over an even number of terms.
\end{proof}
Lemma~\ref{lem:free} uses a different output test (the ciphertext passed through the public map $T^{-1}$) and gives no information about $RK_{r+1}$. We therefore report round counts for $S^{(r)}$ and state the $(r{+}1)$-round extension separately.

\section{Validation}\label{sec:validation}

\paragraph{Reference implementation.} Our Python implementation reproduces both published Dipper-64/128 test vectors and is cross-checked against a vectorised implementation. Section~\ref{sec:erratum} discusses the key schedule.

\paragraph{Local models.} The S-box table is computed from the ANF. We verified the Hu--Yap rule for $x\add y$ exhaustively for $n\le5$, Lemma~\ref{lem:addret} exhaustively for $n\le4$ (all $2^{16}$ exponent pairs), and the carry encoding for $n\le6$.

\paragraph{Benchmark.} With the same SAT machinery and BDP rules we reproduce the results of Eskandari et al.~\cite{eskandari2018} for GIFT-64~\cite{gift2017} and PRESENT~\cite{present2007}: a 9-round property with 63 active bits for GIFT-64 (30 balanced bits when the constant bit is the most significant bit of a nibble, as in the published result; 32 for the best position), no 10-round property, a 9-round property with 60 active bits and one balanced bit for PRESENT, and no 10-round property.

\paragraph{Soundness against experiment.} We evaluated cube sums for random constants and independent random round keys in three places: the word cubes of Table~\ref{tab:words} (1016 trials), the 28 cubes of the tightness study in Section~\ref{sec:tightness} (200 trials per cube and round), and every cube of dimension at most 16 in the frontier of Table~\ref{tab:frontier}, including those of the variants (400 trials). No certified bit ever had a nonzero sum in any of these experiments.

\section{Results on Dipper}\label{sec:results}

\subsection{Reproducing the published experiment}
Table~\ref{tab:words} repeats the published experiment: one 16-bit word saturated, the other 48 bits constant, and 1016 trials (1000 with independent random round keys, 16 with the real 128-bit key schedule). A bit counts as empirically balanced only if its sum is zero in every trial. Five rounds are the maximum, as published, and the round count refers to the state $S^{(r)}$ itself. For every word and every round $r=2,\dots,5$, the exact MP model certifies exactly the empirically balanced bits, so these properties now hold for all keys. The trial count matters: an earlier run with 80 trials reported two additional bits for word $C$ (one at $r=4$, one at $r=5$) whose sums turned out to be nonzero with probability of roughly 0.5--1\%. This is the weakness of purely experimental bounds that certificates remove. In every trial we also confirmed $\bigoplus T^{-1}(S^{(r+1)})=\bigoplus S^{(r)}$ (Lemma~\ref{lem:free}).

\begin{table}[t]\centering
\caption{Balanced bits of $S^{(r)}$ for 16-bit word-saturated cubes. For Dipper ($r=1$--$6$) the empirically zero bits (1016 trials) coincide with the certified bits in every entry for $r\ge2$. The last two columns give, for the analysis variants, the last round with empirically zero bits (400 trials) and their number; these are empirical only. Bold: last round with balanced bits.}\label{tab:words}
\begin{tabular}{lcccccc|cc}
\toprule
Saturated word & $r{=}1$ & 2 & 3 & 4 & 5 & 6 & $\mathrm{Dipper}^{\oplus}$ & $\mathrm{Dipper}^{\varnothing}$\\
\midrule
$A$ & 64 & 64 & 64 & 38 & \textbf{3} & 0 & 6 (26 bits) & 7 (14 bits)\\
$B$ & 64 & 56 & 24 & 0 & 0 & 0 & 5 (7 bits) & 7 (9 bits)\\
$C$ & 64 & 64 & 64 & 38 & \textbf{3} & 0 & 6 (25 bits) & 7 (4 bits)\\
$D$ & 64 & 60 & 25 & 0 & 0 & 0 & 5 (8 bits) & 7 (6 bits)\\
\bottomrule
\end{tabular}
\end{table}

\subsection{Certified properties over all cubes}
By Lemma~\ref{lem:mono}, the 64 cubes of dimension~63 (one constant bit $p$) decide which rounds admit a certificate for some cube. Table~\ref{tab:maximal} summarises them. Six rounds admit certificates for 14 of the 64 positions of the constant bit. The best positions $p\in\{4,6,7\}$ give nine balanced bits $\{1,18,19,24,41,48,51,56,58\}$. At seven rounds no position gives a certificate, so within the model no bit-aligned cube of any dimension has a seven-round certificate for any single output bit. With Lemma~\ref{lem:free}, the six-round properties are seven-round distinguishers on $T^{-1}(C)$.

\begin{table}[t]\centering
\caption{Cubes of dimension 63 (constant bit $p$, 64 cubes per row). Identical for exact MP, MP-circuit and BDP.}\label{tab:maximal}
\begin{tabular}{cccc}
\toprule
Rounds $r$ & cubes with a certificate & max.\ balanced bits & solver time (MP, all 64 cubes) \\
\midrule
4 & 64 & 64 & 3.9 s\\
5 & 64 & 42 & 6.1 s\\
6 & 14 & 9  & 7.9 s\\
7 & 0  & 0  & 8.6 s\\
\bottomrule
\end{tabular}
\end{table}

\paragraph{Data complexity.} Starting from certified cubes, we removed active bits one at a time for as long as some bit stayed certified. We restarted from several cubes and random orders. Table~\ref{tab:frontier} lists the smallest cubes found. They are upper bounds on the data needed, not proven minima. The five-round bound $2^{16}$ coincides with the published word-saturation structure (word $C$). The six-round cube keeps the nibble $\{4,5,6,7\}$ constant and certifies six bits.

\begin{table}[t]\centering
\caption{Smallest certified cubes found (dimension $d$, i.e.\ $2^d$ chosen plaintexts) per round count. ``--'': no certificate for any bit-aligned cube (Lemma~\ref{lem:mono}); blank: not computed.}\label{tab:frontier}
\begin{tabular}{lcccccccc}
\toprule
Rounds $r$ & 3 & 4 & 5 & 6 & 7 & 8 & 9 & 10\\
\midrule
Dipper & 2 & 4 & 16 & 60 & -- & -- & -- & --\\
$\mathrm{Dipper}^{\oplus}$ & & & 7 & 22 & 44 & 59 & 63 & --\\
$\mathrm{Dipper}^{\varnothing}$ & & & 3 & 8 & 15 & 48 & 59 & 63\\
\bottomrule
\end{tabular}
\end{table}

\subsection{Where the balanced bits come from}
Mapping the nine balanced bits at the six-round boundary back through the final bit permutation shows that every one of them is an output of word $B$ or word $D$ in the last round (bits $D_3,D_4,D_5,D_6,D_{11}$ and $B_4,B_{12},B_{13},B_{14}$). These are the words that feed the additions but pass through them unchanged. No output bit of $A\add B$ or $C\add D$ is balanced at the boundary. The same holds for the five-round word-cube properties: bits $\{51,56,58\}$ of the word-$C$ cube map to $D_3,B_4,B_{14}$ and bits $\{1,19,24\}$ of the word-$A$ cube to $D_5,D_{11},B_{12}$. No output bit of the last-round additions is certified at the boundary; only retained-operand bits are. (Bit 0 of $A\add B$ is linear, so this is an observation about the certified bits, not a general statement about every output bit of the additions.)

\subsection{Consequences for key recovery}
Consider an eight-round attack: the six-round property on bit $j$ of $S^{(6)}$, the free seventh round (Lemma~\ref{lem:free}), and one round with a key guess, i.e.\ $S^{(6)}_j\oplus RK_{7,j}=T^{-1}\big(T^{-1}(C)\oplus RK_8\big)_j$. Bit $j$ of $T^{-1}$ depends on four input bits if its nibble lies in word $B$ or $D$ (rotation only), and on up to 32 bits if it lies in $A$ or $C$ (borrow chain of the subtraction). Bit $j=1$ is certified for the $2^{60}$ cube and depends on only four bits of $RK_8$. Each structure of $2^{60}$ plaintexts gives one parity condition on those four bits. We validated the procedure end-to-end on a scaled-down instance (4-round certificate of the cube $\{60,\dots,63\}$ on bit~1, followed by the free round and one key-guessed round, i.e.\ six rounds, 200 random keys). The right guess always survived. A wrong guess survived one structure with probability about $0.565$, so eight structures left on average $1.13$ candidates. Applying this rate to the eight-round attack is an extrapolation from the scaled-down instance. The eight-round attack therefore needs about $8\cdot2^{60}=2^{63}$ chosen plaintexts and $2^{63}$ partial decryptions to recover four bits of $RK_8$. This data complexity is half the codebook, so the attack is of theoretical interest only. Our best integral attack therefore covers 8 of the 28 rounds.

\section{Precision of the models}\label{sec:tightness}

\paragraph{MP versus BDP.} Across all 64 maximal cubes, exact MP and BDP produced identical certificate sets for $r=4,\dots,7$ on Dipper and up to $r=11$ on the variants, and MP-circuit agreed with both wherever it was run ($r\le7$). Across 140 further (cube, round) instances of dimension 4--16 on Dipper, all three models produced \emph{identical} certificate sets. Their costs were similar: 8.4\,s, 8.8\,s and 6.2\,s for the 140 instances. Part of the explanation is structural. In Dipper every S-box layer is preceded by a full-state key addition. In the MP model, the composition ``key addition then S-box'' allows exactly the transitions $u\to v$ with $T[w][v]=1$ for some $w\succeq u$, and this is precisely the BDP S-box table. The two approaches can therefore differ only inside the ARX layer, and there we observed no difference. This agrees with the general observation of Hu et al.~\cite{hu2020mp} that existence-based models lose precision mainly through cancellation, which neither model captures.

\paragraph{Certificates versus experiment.} Table~\ref{tab:tight} compares certificates with experiment on 28 cubes (word, nibble-aligned and random, dimension 4--16). \emph{Empirically zero} means zero in 200 trials. A \emph{gap} bit is empirically zero but not certified. Gap bits were re-tested with 3000 more trials (1000 for 16-dimensional cubes), and \emph{persistent} gaps stayed zero. At four and five rounds the certificates match experiment up to three gap bits, none of them persistent. At two and three rounds about 7--17\% of the empirically zero bits are not certified, and 46 bits stay zero in all re-tests. Such a bit is either an exact property lost to trail cancellation or a sum that is nonzero only with small probability. We attempted to resolve one such bit (bit 30 of a 12-dimensional random cube at two rounds) by exact trail counting with parity per key monomial; the number of trails exceeded our enumeration cap of $10^5$, because every key-addition layer multiplies the trails by key monomials. We leave these cases open. In this sample they occur only at two and three rounds. The frontier itself (six and seven rounds, cubes of dimension 60--63) is out of experimental reach, so we cannot measure tightness there.

\begin{table}[t]\centering
\caption{Certified versus empirically zero bits over 28 cubes of dimension 4--16 (sums over all cubes).}\label{tab:tight}
\begin{tabular}{ccccc}
\toprule
Rounds & certified & empirically zero & gap & persistent gap\\
\midrule
2 & 1283 & 1384 & 101 & 21\\
3 & 446 & 537 & 91 & 25\\
4 & 91 & 93 & 2 & 0\\
5 & 6 & 7 & 1 & 0\\
\bottomrule
\end{tabular}
\end{table}

\section{The role of the modular addition}\label{sec:ablation}

Table~\ref{tab:ablation} compares Dipper with $\mathrm{Dipper}^{\oplus}$ and $\mathrm{Dipper}^{\varnothing}$ under identical cubes (all 64 maximal cubes), models (exact MP and BDP, identical results) and solver settings (conflict budget $10^5$ with a $5\cdot10^6$ retry, no instance left unresolved). Within the models, the two modular additions shorten the longest certified property from nine rounds (XOR) or ten rounds (no addition) to six. The empirical word-cube experiment (Table~\ref{tab:words}) shows the same ordering: 5, 6 and 7 rounds. The data frontier (Table~\ref{tab:frontier}) shifts in the same way: six rounds need $2^{60}$ plaintexts for Dipper, $2^{22}$ for $\mathrm{Dipper}^{\oplus}$ and $2^{8}$ for $\mathrm{Dipper}^{\varnothing}$. The additions are linear in their lowest bit and nonlinear through the carries, whose chains can raise the algebraic degree of the high-order output bits quickly. This is consistent with the observation above that only the retained operands survive at the boundary.

\begin{table}[t]\centering
\caption{Longest round count with a certificate for some cube (all cubes, by Lemma~\ref{lem:mono}) and the maximum number of balanced bits at that round.}\label{tab:ablation}
\begin{tabular}{lccc}
\toprule
Variant & longest certified $r$ & max.\ balanced bits & with free round (Lemma~\ref{lem:free})\\
\midrule
Dipper ($\add$) & 6 & 9 & 7\\
$\mathrm{Dipper}^{\oplus}$ & 9 & 5 & 10\\
$\mathrm{Dipper}^{\varnothing}$ & 10 & 2 & 11\\
\bottomrule
\end{tabular}
\end{table}

A caution about scope: the variants keep the rotations, the S-box, the permutation and the full-state key addition, so they are not GIFT-64. A difference between variants measured with an existence-based model shows how the models behave. It becomes a statement about the true algebraic behaviour only where the models are tight. We measured tightness only on Dipper at up to five rounds (Section~\ref{sec:tightness}), not at the boundary rounds of the variants, so the round differences in Table~\ref{tab:ablation} are statements about certified properties.

\section{Specification check}\label{sec:erratum}
Our text-faithful implementation of Dipper-64/128 reproduces both published test vectors only when the round key is taken as $RK_r=K^{(r)}[63{:}0]$. Equation~(7) of~\cite{dipper2026} states $K^{(r)}[127{:}64]$. For Dipper-64/96 we could not reproduce either published test vector. This remained true under a brute-force search over all 720 word permutations of the key update, all placements and directions of the two word rotations, every 3- or 4-subset of S-box words, all six round-constant positions and three extraction windows. The round function itself is confirmed by the 128-bit vectors. None of the integral results depends on the key schedule, because every certificate holds for arbitrary independent round keys. This is an erratum to our design paper: before submission, the 96-bit schedule must be checked against the reference code, and intermediate round keys should be published with the test vectors.

\section{Related work}
Automated division-property search for ARX ciphers was introduced by Sun, Wang and Wang~\cite{sun2016arx} and extended to SAT~\cite{sun2017sat}. Monomial prediction~\cite{hu2020mp} and 3SDPwoU~\cite{hao2020} give exact algebraic criteria. Hu and Yap~\cite{huyap2024} brought the Braeken--Semaev characterisation~\cite{braeken2005} into monomial prediction for modular addition. CLAASP-MP~\cite{bellini2026claaspmp} is a general MILP framework for monomial prediction that also covers modular addition. We do not claim a new general framework. Our contribution is the analysis of one cipher: exact rules for its particular ARX shape, all-key certificates, a round bound over all cubes within the model, a model comparison and an ablation. Integral key recovery built on monomial prediction is studied, for example, by Hadipour and Eichlseder~\cite{hadipour2022}.

\section{Limitations}
(1) The round bound at seven rounds holds within existence-based models. It does not exclude a seven-round integral property that depends on trail cancellation, or other kinds of integral distinguishers: input sets that are not bit-aligned cubes (affine subspaces in other bases, non-affine sets), or output functions other than single bits. (2) Independent round keys make the certificates stronger (they hold for all keys) but may miss properties that depend on the key schedule. (3) The data-complexity frontier comes from greedy search and gives upper bounds. (4) The gap bits of Section~\ref{sec:tightness} are unresolved. (5) The eight-round key recovery needs $2^{63}$ data and is theoretical.

\section{Conclusion}
The experimental five-round integral bound of Dipper is now a certified all-key result. Certified properties extend to six rounds ($2^{60}$ data), and to seven rounds with the keyless final-round inversion. Within the model no cube reaches seven rounds. In terms of certified properties, the two half-state modular additions cost three to four integral rounds compared with XOR or no addition, and only retained-operand bits are certified at the boundary. On Dipper, exact monomial prediction brings no precision gain over the conventional division property, because a full-state key addition precedes each S-box layer. The remaining imprecision is at two to three rounds and is caused by effects that no existence-based model captures. Dipper's 28 rounds leave a large margin against these properties.

\section*{Data and code availability}
Reference implementation, SAT models, scripts, raw JSON results and this paper's sources: \repo{}. \texttt{make test} runs all validation checks; \texttt{make reproduce} regenerates every table.

\bibliographystyle{plain}
\bibliography{refs}
\end{document}
